\documentclass{article}[13pt]

\usepackage{../uni-notes-eng}
\usepackage{multicol}
\usepackage{graphicx}

\title {Kombinatoryka $\&$ teoria grafów}
\author{by a fish}
\date {21.03.2137}

\begin{document}

\maketitle

\newpage

\begin{center}
    \includegraphics[height=\textheight]{potega.jpg}
\end{center}

\newpage

\section*{SYLABUS - teoria grafów:}
1. Basic concepts: graphs, paths and cycles, complete andbipartite graphs\\
2. Matchings: Hall's Marriage theorem and its variations\\
3. Forbidden subgraphs: complete bipartite and r-partite subgraphs, chromatic numbers, Tur"an's thorem, asymptotic behaviour og edge density, Erd"os-Stone theorem\\
4. Hamiltonian cycles (Dirac's Theorem), Eulerian circuits\\
5. Connectivity: connected and k-connected graphs, Menger's theorem\\
6. Ramsey theory: edge colourings of graphs, Ramsey's theorem and its variations, asymptotic bounds on Ramsey numbers\\
7. Planar graphs and colourings: statements of Kuratowski's and Four Colour theorems, proof of Five Colour theorem, graphs on other surfaces and Euler chracteristics, chromatic polynomial, edge colourings and Vizing's theorem\\
8. Random graphs: further asymptotic bounds on Ramsey numbers, Zarankiewicz numbers and their bounds, graphs of large firth and high chromatic number, cmplete subgraphs in random graphs.\\
9. Algebraic methods: adjavenvy matrix and its eigenvalues, strongly regular graphs, Moore graphs and their existence.

\newpage


\tableofcontents

\newpage

\begin{multicols*}{2}
    
\begin{mdframed}[style=definicja]
    {\color{def}Definicja} - to jest bardzo {\color{acc}wyedukowany} tekst w ramce 
\end{mdframed}

Troszke mniej wyedukowany tekst, bo jest poza fajna ramka ktora wyglada jak guwno i w sumie to nie wiem czemu ja ja w ten sposob ranie
\end{multicols*}

\end{document}